% !TEX program = lualatex
\documentclass{article}

\usepackage{graphicx}
\usepackage{geometry}
\usepackage{xcolor}
\usepackage{hyperref}
\usepackage{enumitem}
\usepackage{caption}
\usepackage{fancyvrb}

\geometry{a4paper, margin=2.5cm}

\hypersetup{colorlinks=true, urlcolor=blue, linkcolor=blue}

\newcommand{\preview}[2][]{%
  \fbox{\includegraphics[#1]{example/#2}}%
}

\title{The \texttt{wrapgraphics} package\\
  \large Wrap text around an image's alpha contour}
\author{}
\date{Version 0.2}

\begin{document}

\maketitle

\begin{abstract}
  \texttt{wrapgraphics} is a Lua\LaTeX{} package that wraps text around an
  image following the contour of its alpha channel. A Python backend
  (Pillow) traces the contour; Lua computes \texttt{\string\parshape} entries;
  the image is overlaid via \texttt{\string\rlap}.
\end{abstract}

\section{Requirements}

\begin{itemize}
  \item Lua\LaTeX{} (with \texttt{--shell-escape})
  \item Python~3 + Pillow
\end{itemize}

\section{Installation}

Place \texttt{wrapgraphics.sty}, \texttt{wrapgraphics.lua}, and
\texttt{wrapgraphics.py} in the same directory as your document
(or in \texttt{\$TEXMF/tex/latex/wrapgraphics/} for system-wide install).

\section{Usage}

\begin{Verbatim}
\usepackage{wrapgraphics}
...
\wrapgraphics[threshold=0.5, padding=5, scale=0.20]{image.png}
\end{Verbatim}

The text following \texttt{\string\wrapgraphics} will flow around the image.
The image is placed as an overlay so it sits inside the paragraph without
breaking the text flow.

\section{Package options}

\begin{description}
  \item[verbose]  Write diagnostic messages to the terminal/log.
  \item[contour]  Draw the traced contour in red over the image
                  (debugging aid).
\end{description}

\section{Key--value reference}

\begin{description}
  \item[threshold] (default \texttt{0.5}) --- Alpha threshold (0--1).
    Pixels with alpha $>$ threshold are considered opaque; the contour
    runs along the boundary of opaque pixels.
  \item[padding] (default \texttt{5}) --- Extra clearance in pixels
    between the image edge and the wrapped text.
  \item[scale] (default \texttt{1}) --- Scale factor passed to
    \texttt{\string\includegraphics}.
  \item[width] --- Explicit width override (e.g.\ \texttt{width=8cm}).
    Overrides \texttt{scale}.
  \item[position] (\texttt{left} or \texttt{right}) --- Which side of
    the image the text wraps around.
  \item[hang] (\texttt{true} or \texttt{false}) --- Use a constant
    indent equal to the widest contour point, rather than following the
    exact shape.
  \item[lines] --- Limit wrapping to the first $N$ lines; remaining
    lines use the full text width.
\end{description}

\section{Examples}

\subsection{Basic usage}

Default \texttt{position=left} with \texttt{threshold=0.5},
\texttt{padding=5}, \texttt{scale=0.20}.

\begin{figure}[h]
  \centering
  \preview[page=1,width=0.75\textwidth]{basic.pdf}
  \caption{Basic left wrapping.}
\end{figure}

\clearpage
\subsection{Left vs.\ right wrapping}

The \texttt{position} key controls which side the text flows around.

\begin{figure}[h]
  \centering
  \preview[page=1,width=0.48\textwidth]{comparison-left.pdf}
  \preview[page=1,width=0.48\textwidth]{comparison-right.pdf}
  \caption{Left (default) and right wrapping.}
\end{figure}

\clearpage
\subsection{Hang mode}

With \texttt{hang=true} the indent is constant for all lines, equal to the
widest contour point.  Text is pushed by a fixed amount rather than
following the exact shape.

\begin{figure}[h]
  \centering
  \preview[page=1,width=0.48\textwidth]{comparison-left-hang.pdf}
  \preview[page=1,width=0.48\textwidth]{comparison-right-hang.pdf}
  \caption{Hang mode (left and right).}
\end{figure}

\clearpage
\subsection{Threshold}

The \texttt{threshold} parameter (0--1) controls how much of the alpha
channel is considered opaque.  A low value treats more pixels as opaque,
keeping the contour closer to the visible image edge.  A high value
lets text come closer to the core of the image.

\begin{figure}[h]
  \centering
  \preview[page=1,width=0.48\textwidth]{params-threshold-03.pdf}
  \preview[page=1,width=0.48\textwidth]{params-threshold-08.pdf}
  \caption{Threshold 0.3 (left) and 0.8 (right).}
\end{figure}

\clearpage
\subsection{Padding}

The \texttt{padding} value (in pixels) adds extra clearance between the
contour and the text.  This is implemented by dilating the alpha mask
before contour tracing.

\begin{figure}[h]
  \centering
  \preview[page=1,width=0.48\textwidth]{params-padding-50.pdf}
  \preview[page=1,width=0.48\textwidth]{params-padding-00.pdf}
  \caption{Padding 50\,px (left) and 0\,px (right).}
\end{figure}

\clearpage
\subsection{Limit lines}

The \texttt{lines} key limits wrapping to the first $N$ lines; subsequent
lines use the full text width.

\begin{figure}[h]
  \centering
  \preview[page=1,width=0.75\textwidth]{lines.pdf}
  \caption{Wrapping limited to 5 lines.}
\end{figure}

\clearpage
\subsection{Width override}

Instead of \texttt{scale}, you can specify an explicit \texttt{width}.

\begin{figure}[h]
  \centering
  \preview[page=1,width=0.75\textwidth]{width.pdf}
  \caption{Explicit width (\texttt{width=8cm}).}
\end{figure}

\clearpage
\subsection{Multiple images}

Multiple \texttt{\string\wrapgraphics} calls can be placed throughout the
document.  Each call independently traces the contour and generates the
appropriate \texttt{\string\parshape}.

\begin{figure}[h]
  \centering
  \preview[page=1,width=0.75\textwidth]{multiple.pdf}
  \caption{Multiple \texttt{\string\wrapgraphics} calls in one document.}
\end{figure}

\clearpage
\subsection{Contour overlay}

With the \texttt{contour} package option, the traced contour is drawn in
red over the image, which is useful for debugging and visualisation.

\begin{figure}[h]
  \centering
  \preview[page=1,width=0.75\textwidth]{verbose.pdf}
  \caption{Contour overlay enabled (\texttt{contour} option).}
\end{figure}

\clearpage
\subsection{Two-column layout}

\texttt{\string\wrapgraphics} works in \texttt{twocolumn} mode.
The wrapping area is sized to the column width.

\begin{figure}[h]
  \centering
  \preview[page=1,width=0.48\textwidth]{twocolumn.pdf}
  \caption{Text wrapping in two-column mode.}
\end{figure}

\clearpage
\section{How it works}

\begin{enumerate}
  \item \texttt{\string\wrapgraphics} saves the image in a box, then
    calls Lua via \texttt{\string\directlua}.
  \item Lua determines the contour SVG path and calls
    \texttt{wrapgraphics.py} (via \texttt{os.execute}) if a cached SVG
    with matching parameters does not already exist.
  \item Python extracts the alpha channel, applies threshold and
    dilation (padding), traces the outer contour with the
    Moore--Neighbor boundary following algorithm, and writes a
    \texttt{.shape.svg} file.
  \item Lua parses the SVG, converts pixel coordinates to
    (\texttt{pt} / \texttt{sp}), and builds a flat list of
    \texttt{\string\parshape} entries (indent + width per line).
  \item The \texttt{\string\parshape} is injected into the TeX stream,
    and the image is placed as an overlay (\texttt{\string\rlap})
    so it sits inside the reflowed paragraph without consuming space.
  \item A \texttt{post\_linebreak\_filter} callback tracks how many
    lines have been consumed and clears the shape when the image height
    is covered, ensuring correct behaviour across page breaks.
\end{enumerate}

\section{Source files}

\begin{description}
  \item[\texttt{wrapgraphics.sty}]  Package file; defines keys, loads Lua.
  \item[\texttt{wrapgraphics.lua}]  Lua module; SVG parsing, parshape
    computation, image placement.
  \item[\texttt{wrapgraphics.py}]   Python CLI; alpha threshold,
    dilation, contour tracing (Moore--Neighbor).
\end{description}

\end{document}
